% Cote 119: offprints, reports, courses, CV and letters (typed and printed part): archive cote 119, site page 215 (PDF page 216).
% typescript. Draft compared: none.
% AI-transcribed by Claude Opus 5.5 (claude-opus-5-5, adaptive thinking, effort high), Message Batches API, batch r119_full (msgbatch_01DwZNm5Y7YsKBWyAN4sFs8w), 2026-11-13; not yet checked by eye.
\scanpage{119/215}
\mspage{. 2 .}

la forme la plus simple (théorème de Minlos). Les
résultats de [6], plus profonds, semblent avoir été
moins bien assimilés par les développements ultérieurs,
mais ils apparaissent comme source d'inspiration dans
un certain nombre de travaux délicats assez récents
sur des inégalités diverses liées à la théorie des
espaces de Banach, notamment ceux de Pelczynsky. Signalons
également les résultats assez fins de [6] et de [8 bis]
sur les propriétés de décroissance de la suite des
valeurs propres de certains opérateurs dans les espaces
de Hilbert et dans les espaces de Banach généraux.

\uline{Références} : L. Schwartz, J. Dieudonné, I. Gelfand,
P. Cartier, J.L. Lions.

2. \uline{Algèbre Homologique} ([8],[9],[19],[9 bis]) .

Depuis 1955, me plaçant au point de vue de
``l'usager'' et non celui du spécialiste, j'ai été amené
continuellement à élargir et à assouplir le langage
de l'algèbre homologique, notamment sous la poussée des
besoins de la géométrie algébrique (théories de dualité,
théories du type Riemann-Roch, cohomologies \hand{$\ell$}-adiques,
cohomologies du type de De Rham, cohomologies cristallines...
Deux directions principales à ces réflexions : développe-%
ment d'une algèbre homologique non commutative (amorcée
dans [10 bis] et systématisée dans la thèse de J. Giraud
[\hand{$*$}]) ; théorie des catégories dérivées (développée
systématiquement par J.L. Verdier, exposée dans
Hartshorne [\hand{$*$}] , Illusie [\hand{$*$}] et [16 SGA 4 Exp. XVIII]) .
Ces deux courants de réflexion sont d'ailleurs loin
d'être épuisés, et sont sans doute appelés à se rejoindre,
soit au sein d'une ``algèbre homogopique'' dont une esquisse
préliminaire a été faite par Quillen [\hand{$*$}], soit dans

\margin{215}
